\chapter{Introduction}

\section{What is a Randomised Algorithm?}

A randomised algorithm is an algorithm that makes random choices during its execution. 
Unlike deterministic algorithms, whose behaviour is completely determined by the input, 
the behaviour of a randomised algorithm may vary between different executions on the same input.\\


\noindent Typically, the algorithm uses a random number generator (\texttt{rng}) to decide between multiple possible actions. 
As a result, quantities such as running time\footnote{Such algorithms are called Las Vegas algorithms.}, 
memory usage, or even correctness\footnote{
Such an algorithm is called a Monte Carlo algorithm. The fact that correctness itself may become probabilistic can 
initially seem undesirable, especially in critical systems where 
absolute correctness is expected. However, in many practical situations, a sufficiently small error probability
is considered acceptable. For example, if an algorithm can be shown to fail with probability at 
most \(10^{-30}\), then the chance of failure may already be smaller than the probability of random hardware faults such as memory
bit flips caused by cosmic radiation, which programmers typically ignore during ordinary software development.
} can become random variables (though usually not all of them at the same time).\\


\noindent The performance of a randomised algorithm is therefore analysed using probability and expectation. 
If $T_{run}$ denotes the runtime of the algorithm, we are usually interested in quantities 
such as $T_{algo}=E[T_{run}]$ over all possible random choices made by the algorithm.\footnote{$T_{algo}$ can still be a 
random variable over the input - so this may \textbf{not} be the same as average case performance 
of the algorithm $T_{avg}$, which is an expected value over all possible inputs.
But for cleanliness of presentation, we will often ignore this distinction and use these terms 
interchangeably or simply use the variable $T$ to mean either of them, depending on context.}

\vspace{0.6em}
Randomised algorithms are often used for the following purposes:
\begin{itemize}
    \item avoid adversarial worst case inputs (common in competitive programming)
    \item simplify algorithm design and implementation
    \item obtain polynomial-time approximate solutions to NP problems (the field of heuristic algorithms)
\end{itemize}

\section{Implementation of Randomness in \texttt{C++}}

Let us make things more concrete with some actual code and implementation details! \texttt{C++} 
provides pseudo-random\footnote{Since computers are deterministic, they rely on parameters like atmospheric noise and temperature for randomness. The details are irrelevant for us... We will take pseudo random as true random for our purposes} number generators through the \texttt{<random>} library. A commonly used generator is \texttt{mt19937}, which is based on the Mersenne Twister algorithm.

\begin{verbatim}
#include <bits/stdc++.h>
mt19937 rng(chrono::steady_clock::now().time_since_epoch().count()); //using time as a seed
\end{verbatim}

This creates an object rng with seed\footnote{Seed values are used in pseudo random generators as a way to recreate the generated random numbers. This can be used for debugging a randomised algorithm. If you don't care about that, 
set it to the current time in milliseconds or something similar and that should be good enough!} as current time which will be passed as a parameter to other functions.
 It has the capability to produce "truly" random 32-bit integers.

\subsection*{Generating Random Integers}

To generate a random integer uniformly from a range \([L,R]\), we use

\begin{verbatim}
uniform_int_distribution<int>uid(L,R);

int x=uid(rng);
int y=uid(rng);
.
.
.
\end{verbatim}

\noindent Each generation takes around constant time.

\subsection*{Generating Random Real Numbers}

To generate a real number uniformly from an interval \([x,y)\), we use

\begin{verbatim}
uniform_real_distribution<double>urd(x,y);

double x=urd(rng);
double y=urd(rng);
.
.
.
\end{verbatim}

\noindent This also takes constant time per generation.

\subsection*{Generating a Random Permutation}

A random permutation of an array (in-place modification) can be generated using \texttt{shuffle}:

\begin{verbatim}
shuffle(a.begin(),a.end(),rng);
\end{verbatim}

\noindent Internally, this does something like:

\begin{verbatim}
for(int i=n-1;i>0;i--)
{
    int j=random integer from [0,i];
    swap(a[i],a[j]);
}
\end{verbatim}

\noindent Deduce why all permutations are equally likely with this method!

\subsection*{Implementing Random Decisions}

Sometimes, an algorithm must choose between multiple actions with specified probabilities. For example, let three actions X,Y,Z be such that they are executed with probabilities:

\[
\Pr(X)=\frac{1}{5},
\qquad
\Pr(Y)=\frac{3}{5},
\qquad
\Pr(Z)=\frac{1}{5}
\]

\noindent One simple way to implement this is to generate a random integer uniformly from \(1\) to \(5\):

\begin{verbatim}
uniform_int_distribution<int>uid(1,5);

int v=uid(rng);

if(v==1)
{
    //do X
}
else if(v<=4)
{
    //do Y
}
else
{
    //do Z
}
\end{verbatim}

thus implementing the required probabilities.


\clearpage
\section{Problems}

\begin{enumerate}

\item[\textbf{1}]
The following procedure is applied on an array \(a\) of \(n\) integers --
\begin{verbatim}
for(int i=n-1;i>0;i--)
{
    int j=random integer from [0,i];
    swap(a[i],a[j]);
}
\end{verbatim}

\begin{enumerate}[label=\textbf{(\alph*)}]
\item Show that this indeed generates a permuation.
\item Show that every permutation is equally likely to be generated.
\end{enumerate}

\item[\textbf{2}]
Suppose you have access only to a function \texttt{bool coin()} which returns \(0\) or \(1\) 
with equal probability. Let the phrase "generate a probability $p$" mean designing a function 
which returns $1$ with probability $p$ and $0$ with probability $1-p$.

\begin{enumerate}[label=\textbf{(\alph*)}]
\item Design a procedure to generate a probability $\frac{k}{2^m}$ using only finite number of calls to \texttt{coin()}.
\item Design a procedure to generate any rational probability $\frac{p}{q}$ with \textbf{expected} 
finite number of calls to \texttt{coin()}\\
\textit{Hint: Think of the binary representation of a number.}
\item Generalise the above ideas to generate any real probability $p$ with \textbf{expected} 
finite number of calls to \texttt{coin()}.
\end{enumerate}

\end{enumerate}